\documentclass[11pt,a4paper]{article}
\usepackage[utf8]{inputenc}
\usepackage[T1]{fontenc}
\usepackage[english]{babel}
\usepackage{geometry}
\usepackage{amsmath}
\usepackage{parskip}

\geometry{top=2.5cm, bottom=2.5cm, left=3cm, right=3cm}

\begin{document}

\begin{center}
  {\LARGE\bfseries GN2: Transformer-Based Jet Flavour Tagging}\\[6pt]
  {\large A PyTorch Implementation}\\[4pt]
  {\small Based on: ATLAS Collaboration,
    \textit{Transforming jet flavour tagging at ATLAS},
    \textit{Nature Communications} \textbf{17}, 541 (2026)}
\end{center}

\vspace{8pt}
\hrule
\vspace{10pt}

\begin{abstract}
Jet flavour tagging is the task of classifying hadronic jets produced in
proton--proton collisions at the Large Hadron Collider according to the flavour
of their originating parton: bottom quark ($b$), charm quark ($c$), hadronic
$\tau$-lepton decay, or light quark/gluon.
This project implements \textbf{GN2}, the transformer-based flavour-tagging
algorithm deployed by the ATLAS Collaboration for Run~2 and Run~3 data
analyses, as a standalone, fully configurable PyTorch pipeline.

The implementation covers the complete machine-learning workflow.
Jet and track variables are read from HDF5 simulation files via a batched,
I/O-optimised data loader that exploits sorted index access for contiguous
reads.
A preprocessing stage applies kinematic selection cuts, performs a
stratified train/val/test split, and computes Z-score normalisation
statistics exclusively on the training set to prevent data leakage.
The GN2 model processes each jet through a per-track initialiser network,
a four-layer transformer encoder with pre-LayerNorm and eight attention heads,
an attention-pooling step that produces a single global jet representation,
and a classification head that outputs four class probabilities
$(p_b, p_c, p_u, p_\tau)$.
Training uses the AdamW optimiser with a linear-warmup and cosine-annealing
learning-rate schedule, reproducing the published ATLAS training recipe.
Evaluation produces standard classification metrics, a normalised confusion
matrix, per-class score distributions, and ROC curves for the $b$- and
$c$-tagging discriminants $D_b$ and $D_c$.

The architecture contains approximately 2.3~million trainable parameters and
does not include the two auxiliary training objectives of the full ATLAS
deployment (track-origin classification and track-pair vertex grouping),
which are identified as the primary target for future extensions.
\end{abstract}

\end{document}